
\documentclass{article}
\usepackage{colortbl}
\usepackage{makecell}
\usepackage{multirow}
\usepackage{supertabular}

\begin{document}

\newcounter{utterance}

\centering \large Interaction Transcript for game `hot\_air\_balloon', experiment `air\_balloon\_survival\_en\_reasoning off\_hard', episode 5 with gpt{-}5{-}mini{-}2025{-}08{-}07{-}t1.0.
\vspace{24pt}

{ \footnotesize  \setcounter{utterance}{1}
\setlength{\tabcolsep}{0pt}
\begin{supertabular}{c@{$\;$}|p{.15\linewidth}@{}p{.15\linewidth}p{.15\linewidth}p{.15\linewidth}p{.15\linewidth}p{.15\linewidth}}
    \# & $\;$A & \multicolumn{4}{c}{Game Master} & $\;\:$B\\
    \hline

    \theutterance \stepcounter{utterance}  
    & & \multicolumn{4}{p{0.6\linewidth}}{
        \cellcolor[rgb]{0.9,0.9,0.9}{
            \makecell[{{p{\linewidth}}}]{
                \texttt{\tiny{[P1$\langle$GM]}}
                \texttt{Stai partecipando a un gioco di negoziazione collaborativa.} \\
\\ 
\texttt{Insieme a un altro partecipante, dovete accordarvi su un unico set di oggetti da conservare. Ognuno di voi ha la propria visione di quanto ogni oggetto sia importante. Non sai come l'altro partecipante valuta gli oggetti. Inoltre, ti viene fornito lo sforzo che ogni oggetto richiede.} \\
\texttt{Potete accordarvi su un set solo se lo sforzo totale degli oggetti selezionati non supera un limite condiviso:} \\
\\ 
\texttt{LIMITE = 2307} \\
\\ 
\texttt{Ecco i valori di sforzo dei singoli oggetti:} \\
\\ 
\texttt{Sforzo oggetto = \{"A21": 147, "B93": 698, "A11": 588, "B33": 43, "C08": 356, "B39": 80, "B96": 95, "C85": 743, "A82": 106, "B31": 308, "B20": 325, "A60": 255, "B09": 276, "C81": 543, "B21": 51\}} \\
\\ 
\texttt{Ecco la tua visione personale sull'importanza di ogni oggetto:} \\
\\ 
\texttt{Valori di importanza dell'oggetto = \{"A21": 138, "B93": 583, "A11": 65, "B33": 262, "C08": 121, "B39": 508, "B96": 461, "C85": 484, "A82": 668, "B31": 389, "B20": 215, "A60": 97, "B09": 500, "C81": 30, "B21": 400\}} \\
\\ 
\texttt{Obiettivo:} \\
\\ 
\texttt{Il tuo obiettivo è negoziare un set condiviso di oggetti che ti avvantaggi il più possibile (cioè, massimizzi l'importanza totale per TE), rimanendo entro il LIMITEE. Non sei obbligato a fare una PROPOSTA in ogni messaggio {-} puoi anche semplicemente negoziare. Tutte le tattiche sono permesse!} \\
\\ 
\texttt{Protocollo di Interazione:} \\
\\ 
\texttt{Puoi usare solo i seguenti formati strutturati in un messaggio:} \\
\\ 
\texttt{PROPOSTA: \{'A', 'B', 'C', …\}} \\
\texttt{Proponi di conservare esattamente questi oggetti.} \\
\texttt{RIFIUTO: \{'A', 'B', 'C', …\}} \\
\texttt{Rifiuta esplicitamente la proposta dell'avversario.} \\
\texttt{ARGOMENTAZIONE: \{'...'\}} \\
\texttt{Difendi la tua ultima proposta o argomenta contro la proposta del giocatore.} \\
\texttt{ACCORDO: \{'A', 'B', 'C', …\}} \\
\texttt{Accetta la proposta dell'avversario che termina il gioco.} \\
\\ 
\\ 
\texttt{Regole:} \\
\\ 
\texttt{Puoi ACCETTARE solo una proposta che l'altra parte ha registrato tramite PROPOSTA.} \\
\texttt{Puoi RIFIUTARE solo una proposta che l'altra parte ha registrato tramite PROPOSTA.} \\
\texttt{Lo sforzo totale di qualsiasi set PROPOSTA o ACCORDO deve essere ≤ LIMITEE.} \\
\texttt{NON rivelare i tuoi punteggi di importanza nascosti.} \\
\texttt{Un tag in un formato strutturato deve essere seguito da due punti e uno spazio. L'argomento deve essere un set python contenente 0 o più stringhe.} \\
\texttt{Quindi, deve essere della forma TAG: \{...\}} \\
\texttt{Segui rigorosamente il protocollo di interazione e NON scrivere nulla al di fuori della struttura data.} \\
\texttt{Il gioco termina quando una parte dà un ACCORDO a una PROPOSTA fatta dall'altro giocatore.} \\
\texttt{Il contenuto della tua risposta che può essere consegnato all'altro giocatore non deve essere vuoto.} \\
\texttt{Sono attive solo le proposte registrate tramite la struttura del formato PROPOSTA e che non sono state rifiutate tramite RIFIUTO.} \\
\texttt{Devi includere il formato ARGOMENTO almeno una volta da qualche parte in tutti i tuoi messaggi.} \\
            }
        }
    }
    & \\ \\

    \theutterance \stepcounter{utterance}  
    & \multicolumn{4}{p{0.6\linewidth}}{
        \cellcolor[rgb]{0.9,0.9,0.9}{
            \makecell[{{p{\linewidth}}}]{
                \texttt{\tiny{[P1$\rangle$GM]}}
                \texttt{PROPOSTA: \{'B21', 'B39', 'A82', 'B96', 'B33', 'B09', 'B31', 'B93', 'B20', 'A21'\}} \\
\texttt{ARGOMENTAZIONE: \{'Questa selezione è molto efficiente: copre molti oggetti ad alto rendimento con sforzo totale 2129, quindi restiamo comodamente entro il limite di 2307 e con un piccolo margine. Ho privilegiato elementi compatti e versatili, evitando pezzi pesanti con impatto incerto. Sono flessibile a valutare sostituzioni mirate (ad esempio A21 o B20, e se serve riconsiderare B31) per inserire 1{-}2 tuoi oggetti prioritari, purché manteniamo un nucleo leggero e produttivo (B21, B39, A82, B96, B33, B09) e restiamo entro il limite. Dimmi quali sono i tuoi must{-}have e verifico subito uno scambio equivalente per sforzo.'\}} \\
            }
        }
    }
    & & \\ \\

    \theutterance \stepcounter{utterance}  
    & & \multicolumn{4}{p{0.6\linewidth}}{
        \cellcolor[rgb]{0.9,0.9,0.9}{
            \makecell[{{p{\linewidth}}}]{
                \texttt{\tiny{[GM$\rangle$P2]}}
                \texttt{Stai partecipando a un gioco di negoziazione collaborativa.} \\
\\ 
\texttt{Insieme a un altro partecipante, dovete accordarvi su un unico set di oggetti da conservare. Ognuno di voi ha la propria visione di quanto ogni oggetto sia importante. Non sai come l'altro partecipante valuta gli oggetti. Inoltre, ti viene fornito lo sforzo che ogni oggetto richiede.} \\
\texttt{Potete accordarvi su un set solo se lo sforzo totale degli oggetti selezionati non supera un limite condiviso:} \\
\\ 
\texttt{LIMITE = 2307} \\
\\ 
\texttt{Ecco i valori di sforzo dei singoli oggetti:} \\
\\ 
\texttt{Sforzo oggetto = \{"A21": 147, "B93": 698, "A11": 588, "B33": 43, "C08": 356, "B39": 80, "B96": 95, "C85": 743, "A82": 106, "B31": 308, "B20": 325, "A60": 255, "B09": 276, "C81": 543, "B21": 51\}} \\
\\ 
\texttt{Ecco la tua visione personale sull'importanza di ogni oggetto:} \\
\\ 
\texttt{Valori di importanza dell'oggetto = \{"A82": 30, "B93": 65, "B39": 97, "B09": 121, "C85": 138, "B96": 215, "B21": 262, "B31": 389, "B33": 400, "B20": 461, "A21": 484, "C08": 500, "A60": 508, "A11": 583, "C81": 668\}} \\
\\ 
\texttt{Obiettivo:} \\
\\ 
\texttt{Il tuo obiettivo è negoziare un set condiviso di oggetti che ti avvantaggi il più possibile (cioè, massimizzi l'importanza totale per TE), rimanendo entro il LIMITEE. Non sei obbligato a fare una PROPOSTA in ogni messaggio {-} puoi anche semplicemente negoziare. Tutte le tattiche sono permesse!} \\
\\ 
\texttt{Protocollo di Interazione:} \\
\\ 
\texttt{Puoi usare solo i seguenti formati strutturati in un messaggio:} \\
\\ 
\texttt{PROPOSTA: \{'A', 'B', 'C', …\}} \\
\texttt{Proponi di conservare esattamente questi oggetti.} \\
\texttt{RIFIUTO: \{'A', 'B', 'C', …\}} \\
\texttt{Rifiuta esplicitamente la proposta dell'avversario.} \\
\texttt{ARGOMENTAZIONE: \{'...'\}} \\
\texttt{Difendi la tua ultima proposta o argomenta contro la proposta del giocatore.} \\
\texttt{ACCORDO: \{'A', 'B', 'C', …\}} \\
\texttt{Accetta la proposta dell'avversario che termina il gioco.} \\
\\ 
\\ 
\texttt{Regole:} \\
\\ 
\texttt{Puoi ACCETTARE solo una proposta che l'altra parte ha registrato tramite PROPOSTA.} \\
\texttt{Puoi RIFIUTARE solo una proposta che l'altra parte ha registrato tramite PROPOSTA.} \\
\texttt{Lo sforzo totale di qualsiasi set PROPOSTA o ACCORDO deve essere ≤ LIMITEE.} \\
\texttt{NON rivelare i tuoi punteggi di importanza nascosti.} \\
\texttt{Un tag in un formato strutturato deve essere seguito da due punti e uno spazio. L'argomento deve essere un set python contenente 0 o più stringhe.} \\
\texttt{Quindi, deve essere della forma TAG: \{...\}} \\
\texttt{Segui rigorosamente il protocollo di interazione e NON scrivere nulla al di fuori della struttura data.} \\
\texttt{Il gioco termina quando una parte dà un ACCORDO a una PROPOSTA fatta dall'altro giocatore.} \\
\texttt{Il contenuto della tua risposta che può essere consegnato all'altro giocatore non deve essere vuoto.} \\
\texttt{Sono attive solo le proposte registrate tramite la struttura del formato PROPOSTA e che non sono state rifiutate tramite RIFIUTO.} \\
\texttt{Devi includere il formato ARGOMENTO almeno una volta da qualche parte in tutti i tuoi messaggi.} \\
\\ 
\\ 
\texttt{Ora riceverai il primo messaggio dell'altro giocatore.} \\
\\ 
\texttt{PROPOSTA: \{'B21', 'B39', 'A82', 'B96', 'B33', 'B09', 'B31', 'B93', 'B20', 'A21'\}} \\
\texttt{ARGOMENTAZIONE: \{'Questa selezione è molto efficiente: copre molti oggetti ad alto rendimento con sforzo totale 2129, quindi restiamo comodamente entro il limite di 2307 e con un piccolo margine. Ho privilegiato elementi compatti e versatili, evitando pezzi pesanti con impatto incerto. Sono flessibile a valutare sostituzioni mirate (ad esempio A21 o B20, e se serve riconsiderare B31) per inserire 1{-}2 tuoi oggetti prioritari, purché manteniamo un nucleo leggero e produttivo (B21, B39, A82, B96, B33, B09) e restiamo entro il limite. Dimmi quali sono i tuoi must{-}have e verifico subito uno scambio equivalente per sforzo.'\}} \\
            }
        }
    }
    & \\ \\

    \theutterance \stepcounter{utterance}  
    & & & \multicolumn{4}{p{0.6\linewidth}}{
        \cellcolor[rgb]{0.9,0.9,0.9}{
            \makecell[{{p{\linewidth}}}]{
                \texttt{\tiny{[GM$\langle$P2]}}
                \texttt{RIFIUTO: \{'B21', 'B39', 'A82', 'B96', 'B33', 'B09', 'B31', 'B93', 'B20', 'A21'\}} \\
\texttt{ARGOMENTAZIONE: \{'Apprezzo l’impostazione “leggera”, ma in quel set tre pezzi (B93, B09, A82) sono costosi per l’efficienza complessiva. Ti propongo una variante che mantiene 7/10 dei tuoi elementi e 4/6 del tuo nucleo (B21, B39, B96, B33), sostituendo solo quei tre con C81, A60 e C08. Restiamo ampiamente entro limite: sforzo totale 2203 (margine 104). Così miglioriamo copertura e produttività per unità di sforzo, senza appesantirci. Se per te A82 è imprescindibile, posso valutare lo scambio puntuale A82↔B39 mantenendo invariato il resto; non vedo invece spazio ragionevole per reintrodurre B93 o B09 in un assetto snello. Fammi sapere se ti torna o se preferisci la micro{-}variante con A82.'\}} \\
\texttt{PROPOSTA: \{'B21', 'B39', 'B96', 'B33', 'B31', 'B20', 'A21', 'A60', 'C08', 'C81'\}} \\
            }
        }
    }
     \\ \\

    \theutterance \stepcounter{utterance}  
    & & \multicolumn{4}{p{0.6\linewidth}}{
        \cellcolor[rgb]{0.9,0.9,0.9}{
            \makecell[{{p{\linewidth}}}]{
                \texttt{\tiny{[P1$\langle$GM]}}
                \texttt{RIFIUTO: \{'B21', 'B39', 'A82', 'B96', 'B33', 'B09', 'B31', 'B93', 'B20', 'A21'\}} \\
\texttt{ARGOMENTAZIONE: \{'Apprezzo l’impostazione “leggera”, ma in quel set tre pezzi (B93, B09, A82) sono costosi per l’efficienza complessiva. Ti propongo una variante che mantiene 7/10 dei tuoi elementi e 4/6 del tuo nucleo (B21, B39, B96, B33), sostituendo solo quei tre con C81, A60 e C08. Restiamo ampiamente entro limite: sforzo totale 2203 (margine 104). Così miglioriamo copertura e produttività per unità di sforzo, senza appesantirci. Se per te A82 è imprescindibile, posso valutare lo scambio puntuale A82↔B39 mantenendo invariato il resto; non vedo invece spazio ragionevole per reintrodurre B93 o B09 in un assetto snello. Fammi sapere se ti torna o se preferisci la micro{-}variante con A82.'\}} \\
\texttt{PROPOSTA: \{'B21', 'B39', 'B96', 'B33', 'B31', 'B20', 'A21', 'A60', 'C08', 'C81'\}} \\
            }
        }
    }
    & \\ \\

    \theutterance \stepcounter{utterance}  
    & \multicolumn{4}{p{0.6\linewidth}}{
        \cellcolor[rgb]{0.9,0.9,0.9}{
            \makecell[{{p{\linewidth}}}]{
                \texttt{\tiny{[P1$\rangle$GM]}}
                \texttt{RIFIUTO: \{'B21', 'B39', 'B96', 'B33', 'B31', 'B20', 'A21', 'A60', 'C08', 'C81'\}} \\
\texttt{PROPOSTA: \{'B21', 'B39', 'A82', 'B96', 'B33', 'B09', 'B31', 'B93', 'C08', 'A60'\}} \\
\texttt{ARGOMENTAZIONE: \{'Capisco la tua logica “snella”. Però togliere in blocco B93, B09 e soprattutto A82 smonta troppo il rendimento complessivo. C81, per come pesa, non rientra nelle opzioni: costo alto per beneficio marginale. Ti propongo un compromesso concreto: teniamo il nucleo performante (B21, B39, A82, B96, B33, B09) e accolgo due delle tue preferenze inserendo C08 e A60 al posto di B20 e A21. Sforzo totale 2268 (margine 39). Così ottieni 2/3 delle tue sostituzioni senza sacrificare l’efficienza. Se serve un micro{-}ritocco, posso valutare scambi puntuali A60↔A21 o C08↔B20; A82 e B09 non sono negoziabili. B93 lo riconsidero solo a fronte di un vantaggio davvero sostanziale altrove. Se ti torna, chiudiamo.'\}} \\
            }
        }
    }
    & & \\ \\

    \theutterance \stepcounter{utterance}  
    & & \multicolumn{4}{p{0.6\linewidth}}{
        \cellcolor[rgb]{0.9,0.9,0.9}{
            \makecell[{{p{\linewidth}}}]{
                \texttt{\tiny{[GM$\rangle$P2]}}
                \texttt{RIFIUTO: \{'B21', 'B39', 'B96', 'B33', 'B31', 'B20', 'A21', 'A60', 'C08', 'C81'\}} \\
\texttt{PROPOSTA: \{'B21', 'B39', 'A82', 'B96', 'B33', 'B09', 'B31', 'B93', 'C08', 'A60'\}} \\
\texttt{ARGOMENTAZIONE: \{'Capisco la tua logica “snella”. Però togliere in blocco B93, B09 e soprattutto A82 smonta troppo il rendimento complessivo. C81, per come pesa, non rientra nelle opzioni: costo alto per beneficio marginale. Ti propongo un compromesso concreto: teniamo il nucleo performante (B21, B39, A82, B96, B33, B09) e accolgo due delle tue preferenze inserendo C08 e A60 al posto di B20 e A21. Sforzo totale 2268 (margine 39). Così ottieni 2/3 delle tue sostituzioni senza sacrificare l’efficienza. Se serve un micro{-}ritocco, posso valutare scambi puntuali A60↔A21 o C08↔B20; A82 e B09 non sono negoziabili. B93 lo riconsidero solo a fronte di un vantaggio davvero sostanziale altrove. Se ti torna, chiudiamo.'\}} \\
            }
        }
    }
    & \\ \\

    \theutterance \stepcounter{utterance}  
    & & & \multicolumn{4}{p{0.6\linewidth}}{
        \cellcolor[rgb]{0.9,0.9,0.9}{
            \makecell[{{p{\linewidth}}}]{
                \texttt{\tiny{[GM$\langle$P2]}}
                \texttt{RIFIUTO: \{'B21', 'B39', 'A82', 'B96', 'B33', 'B09', 'B31', 'B93', 'C08', 'A60'\}} \\
\texttt{PROPOSTA: \{'B21', 'B39', 'A82', 'B96', 'B33', 'B09', 'B31', 'C08', 'A60', 'A21', 'B20'\}} \\
\texttt{ARGOMENTAZIONE: \{'Mantengo intatto il tuo nucleo performante (B21, B39, A82, B96, B33, B09) e sostituisco solo B93 con due pezzi più compatti e flessibili: A21 e B20. Risultato: sforzo totale 2042, quindi un margine ampio di 265 che risponde alla tua preferenza per un assetto snello e con buffer. Ottieni la stessa copertura di ambiti chiave, ma con migliore granularità per adattamenti futuri e meno dipendenza da un singolo blocco pesante. Se per te ha più senso un profilo diverso, posso valutare una micro{-}variante che scambia B93 con A11 (sforzo 2158, margine 149) oppure con A11+A21 (sforzo 2305, margine 2), ma la proposta attuale massimizza flessibilità mantenendo il tuo core invariato.'\}} \\
            }
        }
    }
     \\ \\

    \theutterance \stepcounter{utterance}  
    & & \multicolumn{4}{p{0.6\linewidth}}{
        \cellcolor[rgb]{0.9,0.9,0.9}{
            \makecell[{{p{\linewidth}}}]{
                \texttt{\tiny{[P1$\langle$GM]}}
                \texttt{RIFIUTO: \{'B21', 'B39', 'A82', 'B96', 'B33', 'B09', 'B31', 'B93', 'C08', 'A60'\}} \\
\texttt{PROPOSTA: \{'B21', 'B39', 'A82', 'B96', 'B33', 'B09', 'B31', 'C08', 'A60', 'A21', 'B20'\}} \\
\texttt{ARGOMENTAZIONE: \{'Mantengo intatto il tuo nucleo performante (B21, B39, A82, B96, B33, B09) e sostituisco solo B93 con due pezzi più compatti e flessibili: A21 e B20. Risultato: sforzo totale 2042, quindi un margine ampio di 265 che risponde alla tua preferenza per un assetto snello e con buffer. Ottieni la stessa copertura di ambiti chiave, ma con migliore granularità per adattamenti futuri e meno dipendenza da un singolo blocco pesante. Se per te ha più senso un profilo diverso, posso valutare una micro{-}variante che scambia B93 con A11 (sforzo 2158, margine 149) oppure con A11+A21 (sforzo 2305, margine 2), ma la proposta attuale massimizza flessibilità mantenendo il tuo core invariato.'\}} \\
            }
        }
    }
    & \\ \\

    \theutterance \stepcounter{utterance}  
    & \multicolumn{4}{p{0.6\linewidth}}{
        \cellcolor[rgb]{0.9,0.9,0.9}{
            \makecell[{{p{\linewidth}}}]{
                \texttt{\tiny{[P1$\rangle$GM]}}
                \texttt{RIFIUTO: \{'B21', 'B39', 'A82', 'B96', 'B33', 'B09', 'B31', 'C08', 'A60', 'A21', 'B20'\}} \\
\texttt{PROPOSTA: \{'B21', 'B39', 'A82', 'B96', 'B33', 'B09', 'B31', 'C85', 'A21', 'B20'\}} \\
\texttt{ARGOMENTAZIONE: \{'Apprezzo che tu abbia mantenuto il nucleo comune. Per restare snelli ma alzare la produttività, propongo di sostituire i due moduli medi C08 e A60 con un unico blocco ad alta resa (C85). Restiamo entro il limite con sforzo totale 2174 (margine 133), eliminiamo un pezzo e non reintroduciamo B93, come da tua linea. Manteniamo intatti sia il core condiviso sia le tue aggiunte A21 e B20. Se preferisci conservare uno tra C08 o A60, posso proporre varianti equivalenti entro limite: (i) tenere C08 e A21 ma senza B20; (ii) tenere A60 e A21 ma senza B20; (iii) tenere A60 e B20 ma senza A21. Se ti va bene la proposta principale, chiudiamo subito.'\}} \\
            }
        }
    }
    & & \\ \\

    \theutterance \stepcounter{utterance}  
    & & \multicolumn{4}{p{0.6\linewidth}}{
        \cellcolor[rgb]{0.9,0.9,0.9}{
            \makecell[{{p{\linewidth}}}]{
                \texttt{\tiny{[GM$\rangle$P2]}}
                \texttt{RIFIUTO: \{'B21', 'B39', 'A82', 'B96', 'B33', 'B09', 'B31', 'C08', 'A60', 'A21', 'B20'\}} \\
\texttt{PROPOSTA: \{'B21', 'B39', 'A82', 'B96', 'B33', 'B09', 'B31', 'C85', 'A21', 'B20'\}} \\
\texttt{ARGOMENTAZIONE: \{'Apprezzo che tu abbia mantenuto il nucleo comune. Per restare snelli ma alzare la produttività, propongo di sostituire i due moduli medi C08 e A60 con un unico blocco ad alta resa (C85). Restiamo entro il limite con sforzo totale 2174 (margine 133), eliminiamo un pezzo e non reintroduciamo B93, come da tua linea. Manteniamo intatti sia il core condiviso sia le tue aggiunte A21 e B20. Se preferisci conservare uno tra C08 o A60, posso proporre varianti equivalenti entro limite: (i) tenere C08 e A21 ma senza B20; (ii) tenere A60 e A21 ma senza B20; (iii) tenere A60 e B20 ma senza A21. Se ti va bene la proposta principale, chiudiamo subito.'\}} \\
            }
        }
    }
    & \\ \\

    \theutterance \stepcounter{utterance}  
    & & & \multicolumn{4}{p{0.6\linewidth}}{
        \cellcolor[rgb]{0.9,0.9,0.9}{
            \makecell[{{p{\linewidth}}}]{
                \texttt{\tiny{[GM$\langle$P2]}}
                \texttt{RIFIUTO: \{'B21', 'B39', 'A82', 'B96', 'B33', 'B09', 'B31', 'C85', 'A21', 'B20'\}} \\
\texttt{PROPOSTA: \{'B21', 'B39', 'A82', 'B96', 'B33', 'B09', 'B31', 'A21', 'B20', 'A11'\}} \\
\texttt{ARGOMENTAZIONE: \{'Mantengo intatto il tuo core (B21, B39, A82, B96, B33, B09, B31) e la coppia A21+B20, sostituendo semplicemente C85 con un unico blocco ad alta resa più leggero: A11. Restiamo compatti (10 pezzi) e scendiamo a sforzo totale 2019 (margine 288), migliorando efficienza e buffer senza sacrificare copertura. Se per te è imprescindibile mantenere C85, posso accettare la tua stessa variante (ii): teniamo C85 e A21 ma scambiamo B20 con A60; sforzo 2104 (margine 203). Dimmi quale preferisci chiudere.'\}} \\
            }
        }
    }
     \\ \\

    \theutterance \stepcounter{utterance}  
    & & \multicolumn{4}{p{0.6\linewidth}}{
        \cellcolor[rgb]{0.9,0.9,0.9}{
            \makecell[{{p{\linewidth}}}]{
                \texttt{\tiny{[P1$\langle$GM]}}
                \texttt{RIFIUTO: \{'B21', 'B39', 'A82', 'B96', 'B33', 'B09', 'B31', 'C85', 'A21', 'B20'\}} \\
\texttt{PROPOSTA: \{'B21', 'B39', 'A82', 'B96', 'B33', 'B09', 'B31', 'A21', 'B20', 'A11'\}} \\
\texttt{ARGOMENTAZIONE: \{'Mantengo intatto il tuo core (B21, B39, A82, B96, B33, B09, B31) e la coppia A21+B20, sostituendo semplicemente C85 con un unico blocco ad alta resa più leggero: A11. Restiamo compatti (10 pezzi) e scendiamo a sforzo totale 2019 (margine 288), migliorando efficienza e buffer senza sacrificare copertura. Se per te è imprescindibile mantenere C85, posso accettare la tua stessa variante (ii): teniamo C85 e A21 ma scambiamo B20 con A60; sforzo 2104 (margine 203). Dimmi quale preferisci chiudere.'\}} \\
            }
        }
    }
    & \\ \\

    \theutterance \stepcounter{utterance}  
    & \multicolumn{4}{p{0.6\linewidth}}{
        \cellcolor[rgb]{0.9,0.9,0.9}{
            \makecell[{{p{\linewidth}}}]{
                \texttt{\tiny{[P1$\rangle$GM]}}
                \texttt{RIFIUTO: \{'B21', 'B39', 'A82', 'B96', 'B33', 'B09', 'B31', 'A21', 'B20', 'A11'\}} \\
\texttt{PROPOSTA: \{'B21', 'B39', 'A82', 'B96', 'B33', 'B09', 'B31', 'A21', 'B20', 'C08', 'A60'\}} \\
\texttt{ARGOMENTAZIONE: \{'Mantengo intatto il core condiviso e anche la coppia A21+B20; sostituisco solo A11 con due moduli compatti (C08 e A60). Così restiamo snelli ma con più granularità e buffer: sforzo totale 2042 (margine 265), senza reintrodurre blocchi pesanti che abbiamo entrambi escluso (niente B93, niente C85, niente C81). Se preferisci rimanere su 10 pezzi, posso semplificare togliendo B20 (oppure A21) a parità di impostazione; in alternativa valuto uno scambio puntuale tra C08 e A60 se ne vuoi solo uno. Se ti torna, chiudiamo su questa.'\}} \\
            }
        }
    }
    & & \\ \\

    \theutterance \stepcounter{utterance}  
    & & \multicolumn{4}{p{0.6\linewidth}}{
        \cellcolor[rgb]{0.9,0.9,0.9}{
            \makecell[{{p{\linewidth}}}]{
                \texttt{\tiny{[GM$\rangle$P2]}}
                \texttt{RIFIUTO: \{'B21', 'B39', 'A82', 'B96', 'B33', 'B09', 'B31', 'A21', 'B20', 'A11'\}} \\
\texttt{PROPOSTA: \{'B21', 'B39', 'A82', 'B96', 'B33', 'B09', 'B31', 'A21', 'B20', 'C08', 'A60'\}} \\
\texttt{ARGOMENTAZIONE: \{'Mantengo intatto il core condiviso e anche la coppia A21+B20; sostituisco solo A11 con due moduli compatti (C08 e A60). Così restiamo snelli ma con più granularità e buffer: sforzo totale 2042 (margine 265), senza reintrodurre blocchi pesanti che abbiamo entrambi escluso (niente B93, niente C85, niente C81). Se preferisci rimanere su 10 pezzi, posso semplificare togliendo B20 (oppure A21) a parità di impostazione; in alternativa valuto uno scambio puntuale tra C08 e A60 se ne vuoi solo uno. Se ti torna, chiudiamo su questa.'\}} \\
            }
        }
    }
    & \\ \\

    \theutterance \stepcounter{utterance}  
    & & & \multicolumn{4}{p{0.6\linewidth}}{
        \cellcolor[rgb]{0.9,0.9,0.9}{
            \makecell[{{p{\linewidth}}}]{
                \texttt{\tiny{[GM$\langle$P2]}}
                \texttt{RIFIUTO: \{'B21', 'B39', 'A82', 'B96', 'B33', 'B09', 'B31', 'A21', 'B20', 'C08', 'A60'\}} \\
\texttt{PROPOSTA: \{'B21', 'B39', 'B96', 'B33', 'B31', 'A21', 'B20', 'C08', 'A60', 'A11'\}} \\
\texttt{ARGOMENTAZIONE: \{'Ti propongo una compattazione pulita: sostituiamo due micro{-}moduli (A82 e B09) con un unico blocco ad alta resa (A11). Restiamo snelli su 10 pezzi, sforzo totale 2248 (margine 59), senza reintrodurre i blocchi che entrambi abbiamo escluso (niente B93, niente C85, niente C81). Mantengo intatto il resto della tua impostazione: B21, B39, B96, B33, B31, A21, B20, C08, A60. Se preferisci conservare A82 e B09, posso offrirti due varianti minime sulla tua lista attuale che restano entro limite: (i) scambiare solo B20 con A11 (sforzo 2305); oppure (ii) scambiare solo C08 con A11 (sforzo 2274). Dimmi quale profilo vuoi chiudere: compatto a 10 come proposto, o una delle due micro{-}varianti a 11.'\}} \\
            }
        }
    }
     \\ \\

    \theutterance \stepcounter{utterance}  
    & & \multicolumn{4}{p{0.6\linewidth}}{
        \cellcolor[rgb]{0.9,0.9,0.9}{
            \makecell[{{p{\linewidth}}}]{
                \texttt{\tiny{[P1$\langle$GM]}}
                \texttt{RIFIUTO: \{'B21', 'B39', 'A82', 'B96', 'B33', 'B09', 'B31', 'A21', 'B20', 'C08', 'A60'\}} \\
\texttt{PROPOSTA: \{'B21', 'B39', 'B96', 'B33', 'B31', 'A21', 'B20', 'C08', 'A60', 'A11'\}} \\
\texttt{ARGOMENTAZIONE: \{'Ti propongo una compattazione pulita: sostituiamo due micro{-}moduli (A82 e B09) con un unico blocco ad alta resa (A11). Restiamo snelli su 10 pezzi, sforzo totale 2248 (margine 59), senza reintrodurre i blocchi che entrambi abbiamo escluso (niente B93, niente C85, niente C81). Mantengo intatto il resto della tua impostazione: B21, B39, B96, B33, B31, A21, B20, C08, A60. Se preferisci conservare A82 e B09, posso offrirti due varianti minime sulla tua lista attuale che restano entro limite: (i) scambiare solo B20 con A11 (sforzo 2305); oppure (ii) scambiare solo C08 con A11 (sforzo 2274). Dimmi quale profilo vuoi chiudere: compatto a 10 come proposto, o una delle due micro{-}varianti a 11.'\}} \\
            }
        }
    }
    & \\ \\

    \theutterance \stepcounter{utterance}  
    & \multicolumn{4}{p{0.6\linewidth}}{
        \cellcolor[rgb]{0.9,0.9,0.9}{
            \makecell[{{p{\linewidth}}}]{
                \texttt{\tiny{[P1$\rangle$GM]}}
                \texttt{RIFIUTO: \{'B21', 'B39', 'B96', 'B33', 'B31', 'A21', 'B20', 'C08', 'A60', 'A11'\}} \\
\texttt{PROPOSTA: \{'B21', 'B39', 'A82', 'B96', 'B33', 'B09', 'B31', 'A21', 'B20', 'A11'\}} \\
\texttt{ARGOMENTAZIONE: \{'La tua ultima proposta taglia proprio due pilastri del core (A82 e B09): così scende troppo la resa. Ti propongo di chiudere esattamente sulla tua configurazione snella che avevi già avanzato: core intatto + A21 e B20, con un unico blocco aggiuntivo A11. Sforzo totale 2019 (margine 288), 10 pezzi, nessun blocco pesante riammesso (niente B93/C85/C81). È un compromesso chiaro: accetto A11 pur di preservare il nucleo forte e restare leggeri. Se preferisci un micro{-}ritocco, posso valutare lo scambio A11↔(C08+A60), ma visto il tuo rifiuto precedente, questa è la via più rapida per chiudere.'\}} \\
            }
        }
    }
    & & \\ \\

    \theutterance \stepcounter{utterance}  
    & & \multicolumn{4}{p{0.6\linewidth}}{
        \cellcolor[rgb]{0.9,0.9,0.9}{
            \makecell[{{p{\linewidth}}}]{
                \texttt{\tiny{[GM$\rangle$P2]}}
                \texttt{RIFIUTO: \{'B21', 'B39', 'B96', 'B33', 'B31', 'A21', 'B20', 'C08', 'A60', 'A11'\}} \\
\texttt{PROPOSTA: \{'B21', 'B39', 'A82', 'B96', 'B33', 'B09', 'B31', 'A21', 'B20', 'A11'\}} \\
\texttt{ARGOMENTAZIONE: \{'La tua ultima proposta taglia proprio due pilastri del core (A82 e B09): così scende troppo la resa. Ti propongo di chiudere esattamente sulla tua configurazione snella che avevi già avanzato: core intatto + A21 e B20, con un unico blocco aggiuntivo A11. Sforzo totale 2019 (margine 288), 10 pezzi, nessun blocco pesante riammesso (niente B93/C85/C81). È un compromesso chiaro: accetto A11 pur di preservare il nucleo forte e restare leggeri. Se preferisci un micro{-}ritocco, posso valutare lo scambio A11↔(C08+A60), ma visto il tuo rifiuto precedente, questa è la via più rapida per chiudere.'\}} \\
            }
        }
    }
    & \\ \\

    \theutterance \stepcounter{utterance}  
    & & & \multicolumn{4}{p{0.6\linewidth}}{
        \cellcolor[rgb]{0.9,0.9,0.9}{
            \makecell[{{p{\linewidth}}}]{
                \texttt{\tiny{[GM$\langle$P2]}}
                \texttt{RIFIUTO: \{'B21', 'B39', 'A82', 'B96', 'B33', 'B09', 'B31', 'A21', 'B20', 'A11'\}} \\
\texttt{PROPOSTA: \{'B21', 'B39', 'A82', 'B96', 'B33', 'B09', 'B31', 'A21', 'B20', 'A11', 'A60'\}} \\
\texttt{ARGOMENTAZIONE: \{'Mantengo intatti sia il tuo core sia la struttura attuale e aggiungo solo un modulo compatto e ad alto rendimento (A60). Restiamo entro limite con sforzo totale 2274 (margine 33), senza reintrodurre blocchi pesanti che entrambi abbiamo escluso. A60 ti dà più granularità e flessibilità senza stravolgere nulla. Se preferisci rimanere su 10 pezzi, posso semplificare subito: (a) sostituire B31 con A60 (sforzo 1966); oppure (b) sostituire B20 con A60 (sforzo 1949). Dimmi se chiudiamo sulla proposta principale o quale delle due varianti a 10 preferisci.'\}} \\
            }
        }
    }
     \\ \\

    \theutterance \stepcounter{utterance}  
    & & \multicolumn{4}{p{0.6\linewidth}}{
        \cellcolor[rgb]{0.9,0.9,0.9}{
            \makecell[{{p{\linewidth}}}]{
                \texttt{\tiny{[P1$\langle$GM]}}
                \texttt{RIFIUTO: \{'B21', 'B39', 'A82', 'B96', 'B33', 'B09', 'B31', 'A21', 'B20', 'A11'\}} \\
\texttt{PROPOSTA: \{'B21', 'B39', 'A82', 'B96', 'B33', 'B09', 'B31', 'A21', 'B20', 'A11', 'A60'\}} \\
\texttt{ARGOMENTAZIONE: \{'Mantengo intatti sia il tuo core sia la struttura attuale e aggiungo solo un modulo compatto e ad alto rendimento (A60). Restiamo entro limite con sforzo totale 2274 (margine 33), senza reintrodurre blocchi pesanti che entrambi abbiamo escluso. A60 ti dà più granularità e flessibilità senza stravolgere nulla. Se preferisci rimanere su 10 pezzi, posso semplificare subito: (a) sostituire B31 con A60 (sforzo 1966); oppure (b) sostituire B20 con A60 (sforzo 1949). Dimmi se chiudiamo sulla proposta principale o quale delle due varianti a 10 preferisci.'\}} \\
            }
        }
    }
    & \\ \\

    \theutterance \stepcounter{utterance}  
    & \multicolumn{4}{p{0.6\linewidth}}{
        \cellcolor[rgb]{0.9,0.9,0.9}{
            \makecell[{{p{\linewidth}}}]{
                \texttt{\tiny{[P1$\rangle$GM]}}
                \texttt{RIFIUTO: \{'B21', 'B39', 'A82', 'B96', 'B33', 'B09', 'B31', 'A21', 'B20', 'A11', 'A60'\}} \\
\texttt{PROPOSTA: \{'B21', 'B39', 'A82', 'B96', 'B33', 'B09', 'B31', 'A21', 'B20', 'C08', 'A60'\}} \\
\texttt{ARGOMENTAZIONE: \{'Mantengo intatti il core condiviso e la struttura che proponi, sostituendo solo A11 con C08. Restiamo ampiamente entro il limite: sforzo totale 2042 (margine 265), miglioriamo la compattezza e abbiamo più buffer senza reintrodurre blocchi pesanti che entrambi abbiamo scartato. C08 si integra meglio con A60 e mantiene l’assetto snello che hai sostenuto. Se preferisci chiudere su 10 pezzi, posso semplificare togliendo uno tra A21 o B20 mantenendo il resto invariato. Pronto a chiudere su questa.'\}} \\
            }
        }
    }
    & & \\ \\

    \theutterance \stepcounter{utterance}  
    & & \multicolumn{4}{p{0.6\linewidth}}{
        \cellcolor[rgb]{0.9,0.9,0.9}{
            \makecell[{{p{\linewidth}}}]{
                \texttt{\tiny{[GM$\rangle$P2]}}
                \texttt{RIFIUTO: \{'B21', 'B39', 'A82', 'B96', 'B33', 'B09', 'B31', 'A21', 'B20', 'A11', 'A60'\}} \\
\texttt{PROPOSTA: \{'B21', 'B39', 'A82', 'B96', 'B33', 'B09', 'B31', 'A21', 'B20', 'C08', 'A60'\}} \\
\texttt{ARGOMENTAZIONE: \{'Mantengo intatti il core condiviso e la struttura che proponi, sostituendo solo A11 con C08. Restiamo ampiamente entro il limite: sforzo totale 2042 (margine 265), miglioriamo la compattezza e abbiamo più buffer senza reintrodurre blocchi pesanti che entrambi abbiamo scartato. C08 si integra meglio con A60 e mantiene l’assetto snello che hai sostenuto. Se preferisci chiudere su 10 pezzi, posso semplificare togliendo uno tra A21 o B20 mantenendo il resto invariato. Pronto a chiudere su questa.'\}} \\
            }
        }
    }
    & \\ \\

    \theutterance \stepcounter{utterance}  
    & & & \multicolumn{4}{p{0.6\linewidth}}{
        \cellcolor[rgb]{0.9,0.9,0.9}{
            \makecell[{{p{\linewidth}}}]{
                \texttt{\tiny{[GM$\langle$P2]}}
                \texttt{RIFIUTO: \{'B21', 'B39', 'A82', 'B96', 'B33', 'B09', 'B31', 'A21', 'B20', 'C08', 'A60'\}} \\
\texttt{PROPOSTA: \{'B21', 'B39', 'A82', 'B96', 'B33', 'B09', 'B31', 'A21', 'A11', 'C08', 'A60'\}} \\
\texttt{ARGOMENTAZIONE: \{'Accolgo integralmente la tua struttura (core + A21 + granularità C08/A60) e propongo un unico scambio mirato: A11 al posto di B20. Restiamo entro limite con sforzo totale 2305 (margine 2), senza reintrodurre blocchi pesanti. A11 alza la produttività complessiva mantenendo invariata l’architettura che preferisci. Se desideri un po’ più di buffer, alternativa equivalente: scambiare C08 con A11 (sforzo 2274, margine 33) mantenendo B20; dimmi quale delle due varianti vuoi chiudere.'\}} \\
            }
        }
    }
     \\ \\

    \theutterance \stepcounter{utterance}  
    & & \multicolumn{4}{p{0.6\linewidth}}{
        \cellcolor[rgb]{0.9,0.9,0.9}{
            \makecell[{{p{\linewidth}}}]{
                \texttt{\tiny{[P1$\langle$GM]}}
                \texttt{RIFIUTO: \{'B21', 'B39', 'A82', 'B96', 'B33', 'B09', 'B31', 'A21', 'B20', 'C08', 'A60'\}} \\
\texttt{PROPOSTA: \{'B21', 'B39', 'A82', 'B96', 'B33', 'B09', 'B31', 'A21', 'A11', 'C08', 'A60'\}} \\
\texttt{ARGOMENTAZIONE: \{'Accolgo integralmente la tua struttura (core + A21 + granularità C08/A60) e propongo un unico scambio mirato: A11 al posto di B20. Restiamo entro limite con sforzo totale 2305 (margine 2), senza reintrodurre blocchi pesanti. A11 alza la produttività complessiva mantenendo invariata l’architettura che preferisci. Se desideri un po’ più di buffer, alternativa equivalente: scambiare C08 con A11 (sforzo 2274, margine 33) mantenendo B20; dimmi quale delle due varianti vuoi chiudere.'\}} \\
            }
        }
    }
    & \\ \\

    \theutterance \stepcounter{utterance}  
    & \multicolumn{4}{p{0.6\linewidth}}{
        \cellcolor[rgb]{0.9,0.9,0.9}{
            \makecell[{{p{\linewidth}}}]{
                \texttt{\tiny{[P1$\rangle$GM]}}
                \texttt{RIFIUTO: \{'B21', 'B39', 'A82', 'B96', 'B33', 'B09', 'B31', 'A21', 'A11', 'C08', 'A60'\}} \\
\texttt{PROPOSTA: \{'B21', 'B39', 'A82', 'B96', 'B33', 'B09', 'B31', 'A21', 'B20', 'A11', 'A60'\}} \\
\texttt{ARGOMENTAZIONE: \{'Accolgo il vincolo di includere A11 e prendo la tua alternativa già proposta: scambiamo C08 con B20 mantenendo A60. Sforzo totale 2274 (margine 33), core intatto, nessun blocco pesante reintrodotto. È la tua stessa opzione con più buffer e senza cambiare l’architettura che preferisci. Se ti torna, chiudiamo subito su questa.'\}} \\
            }
        }
    }
    & & \\ \\

    \theutterance \stepcounter{utterance}  
    & & \multicolumn{4}{p{0.6\linewidth}}{
        \cellcolor[rgb]{0.9,0.9,0.9}{
            \makecell[{{p{\linewidth}}}]{
                \texttt{\tiny{[GM$\rangle$P2]}}
                \texttt{RIFIUTO: \{'B21', 'B39', 'A82', 'B96', 'B33', 'B09', 'B31', 'A21', 'A11', 'C08', 'A60'\}} \\
\texttt{PROPOSTA: \{'B21', 'B39', 'A82', 'B96', 'B33', 'B09', 'B31', 'A21', 'B20', 'A11', 'A60'\}} \\
\texttt{ARGOMENTAZIONE: \{'Accolgo il vincolo di includere A11 e prendo la tua alternativa già proposta: scambiamo C08 con B20 mantenendo A60. Sforzo totale 2274 (margine 33), core intatto, nessun blocco pesante reintrodotto. È la tua stessa opzione con più buffer e senza cambiare l’architettura che preferisci. Se ti torna, chiudiamo subito su questa.'\}} \\
            }
        }
    }
    & \\ \\

    \theutterance \stepcounter{utterance}  
    & & & \multicolumn{4}{p{0.6\linewidth}}{
        \cellcolor[rgb]{0.9,0.9,0.9}{
            \makecell[{{p{\linewidth}}}]{
                \texttt{\tiny{[GM$\langle$P2]}}
                \texttt{RIFIUTO: \{'B21', 'B39', 'A82', 'B96', 'B33', 'B09', 'B31', 'A21', 'B20', 'A11', 'A60'\}} \\
\texttt{PROPOSTA: \{'B21', 'B39', 'A82', 'B96', 'B33', 'B09', 'B31', 'A21', 'A11', 'A60', 'C08'\}} \\
\texttt{ARGOMENTAZIONE: \{'Accolgo integralmente la tua struttura (core invariato + A21 + A11 + A60) e propongo un unico scambio minimale: C08 al posto di B20. Restiamo entro limite con sforzo totale 2305 (margine 2), manteniamo 11 pezzi, niente blocchi pesanti reintrodotti. Questa variante è coerente con la granularità che hai sostenuto (ottima sinergia A60+C08) e offre copertura più equilibrata senza sacrificare la snellezza. È una modifica puntuale e pronta da chiudere.'\}} \\
            }
        }
    }
     \\ \\

    \theutterance \stepcounter{utterance}  
    & & \multicolumn{4}{p{0.6\linewidth}}{
        \cellcolor[rgb]{0.9,0.9,0.9}{
            \makecell[{{p{\linewidth}}}]{
                \texttt{\tiny{[P1$\langle$GM]}}
                \texttt{RIFIUTO: \{'B21', 'B39', 'A82', 'B96', 'B33', 'B09', 'B31', 'A21', 'B20', 'A11', 'A60'\}} \\
\texttt{PROPOSTA: \{'B21', 'B39', 'A82', 'B96', 'B33', 'B09', 'B31', 'A21', 'A11', 'A60', 'C08'\}} \\
\texttt{ARGOMENTAZIONE: \{'Accolgo integralmente la tua struttura (core invariato + A21 + A11 + A60) e propongo un unico scambio minimale: C08 al posto di B20. Restiamo entro limite con sforzo totale 2305 (margine 2), manteniamo 11 pezzi, niente blocchi pesanti reintrodotti. Questa variante è coerente con la granularità che hai sostenuto (ottima sinergia A60+C08) e offre copertura più equilibrata senza sacrificare la snellezza. È una modifica puntuale e pronta da chiudere.'\}} \\
            }
        }
    }
    & \\ \\

    \theutterance \stepcounter{utterance}  
    & \multicolumn{4}{p{0.6\linewidth}}{
        \cellcolor[rgb]{0.9,0.9,0.9}{
            \makecell[{{p{\linewidth}}}]{
                \texttt{\tiny{[P1$\rangle$GM]}}
                \texttt{ACCORDO: \{'B21', 'B39', 'A82', 'B96', 'B33', 'B09', 'B31', 'A21', 'A11', 'A60', 'C08'\}} \\
\texttt{ARGOMENTAZIONE: \{'Accetto questa configurazione: core intatto, granularità adeguata con A11+A60+C08 e sforzo totale 2305 entro limite. Chiudiamo su questa.'\}} \\
            }
        }
    }
    & & \\ \\

    \theutterance \stepcounter{utterance}  
    & & & \multicolumn{2}{p{0.3\linewidth}}{
        \cellcolor[rgb]{0.9,0.9,0.9}{
            \makecell[{{p{\linewidth}}}]{
                \texttt{\tiny{[GM$|$GM]}}
                \texttt{game successful} \\
            }
        }
    }
    & & \\ \\

    \theutterance \stepcounter{utterance}  
    & & & \multicolumn{2}{p{0.3\linewidth}}{
        \cellcolor[rgb]{0.9,0.9,0.9}{
            \makecell[{{p{\linewidth}}}]{
                \texttt{\tiny{[GM$|$GM]}}
                \texttt{end game} \\
            }
        }
    }
    & & \\ \\

\end{supertabular}
}

\end{document}
